% ============================================
% DESEMPENHO DO AM CONVENCIONAL
% Mesmo roteiro do DSB-SC: ponto 1 antes do detector de envoltória,
% ponto 2 depois do bloqueio de CC. No fim, o limiar com ruído forte.
% ============================================

% Diagrama compacto do receptor AM. Argumento: ponto destacado, 1 ou 2.
\newcommand{\diagAM}[1]{%
\begin{tikzpicture}[font=\scriptsize,
  blk/.style={rectangle, draw, rounded corners, minimum height=0.6cm, minimum width=1.2cm, align=center},
  sm/.style={circle, draw, inner sep=0.5pt}]
\node (in) {$s(t)$};
\node[sm, right=0.5cm of in] (sum) {$+$};
\node[above=0.25cm of sum, cRed] (n) {$n(t)$};
\node[blk, right=0.5cm of sum] (bpf) {PF $2W$\\ em $f_c$};
\node[blk, right=1.1cm of bpf] (env) {Detector de\\ envoltória};
\node[blk, right=0.5cm of env] (dc) {Bloqueio\\ de CC};
\node[right=1.1cm of dc] (out) {$y(t)$};
\draw[->] (in) -- (sum);
\draw[->, cRed] (n) -- (sum);
\draw[->] (sum) -- (bpf);
\draw[->] (bpf) -- (env);
\draw[->] (env) -- (dc);
\draw[->] (dc) -- (out);
\coordinate (p1) at ($(bpf.east)!0.5!(env.west)$);
\coordinate (p2) at ($(dc.east)!0.5!(out.west)$);
\ifnum#1=1
  \node[marca, fill=cBlue, text=white] at (p1) {1};
\else
  \node[marca, draw=gray!60, fill=white, text=gray] at (p1) {1};
\fi
\ifnum#1=2
  \node[marca, fill=cGreen, text=white] at (p2) {2};
\else
  \node[marca, draw=gray!60, fill=white, text=gray] at (p2) {2};
\fi
\end{tikzpicture}}

\newcommand{\topoAM}[1]{{\centering\diagAM{#1}\par}\vspace{2pt}}

\subsection{AM convencional}

\begin{frame}{Receptor AM convencional}

\[
s(t) = \big[\tikzmarknode{amA}{A} + m(t)\big]\cos(2\pi f_c t), \qquad A + m(t) \ge 0
\]
\anotacoes{\only<2->{
  \node[nota=cBlue, anchor=east] (a) at ([xshift=-1.6cm,yshift=-0.1cm]amA.west) {portadora transmitida};
  \draw[setinha, cBlue] (a.east) -- ([xshift=-0.3cm]amA.west);
}}%

\begin{center}
\begin{tikzpicture}
\begin{axis}[onda, width=10cm, height=1.5cm, xmin=0, xmax=10, ymin=-1.8, ymax=1.8,
             xlabel={$t$}, ylabel={$s(t)$}]
  \addplot[cBlue, domain=0:10, samples=900] {(1+0.6*cos(deg(2*pi*0.2*x)))*cos(deg(2*pi*3*x))};
  \only<3->{
    \addplot[cRed, thick, dashed, domain=0:10, samples=200] {1+0.6*cos(deg(2*pi*0.2*x))};
    \node[cRed, font=\scriptsize, anchor=west] at (axis cs:10.1,1.6) {$A + m(t)$};
  }
\end{axis}
\end{tikzpicture}
\end{center}

\vspace{0.1cm}
\small
\only<1>{A diferença para o DSB-SC é a parcela $A\cos(2\pi f_c t)$.}%
\only<2>{Com $A + m(t) \ge 0$, a envoltória nunca troca de sinal.}%
\only<3>{O detector de envoltória segue a curva tracejada e dispensa a portadora local.}%
\only<4->{Mesmo roteiro do DSB-SC: sinal, ruído e SNR no ponto 1, antes do detector, e no ponto 2, depois do bloqueio de CC.

\vspace{0.1cm}
{\centering\diagAM{0}\par}}

\end{frame}

% ============================================

\begin{frame}{AM, ponto 1: potência do sinal}

\small\compacto
Mesma conta do DSB-SC, agora com $A + m$ no lugar de $A_c m$:
\onslide<2->{
\[
S_i = \overline{\big[A + m\big]^2\tikzmarknode{c2}{\cos^2(2\pi f_c t)}}
\only<3->{= \tfrac12\,\overline{\big[A + m\big]^2}}
\only<4->{= \tfrac12\big(A^2 + 2A\,\tikzmarknode{mm}{\overline{m}} + \overline{m^2}\big)}
\]}%
\par\vspace{0.7cm}%
\anotacoes{
  \only<3->{
    \node[nota=cBlue, anchor=north] (a) at ([yshift=-0.25cm]c2.south) {média $\tfrac12$: $\cos 4\pi f_c t$ tem média nula};
    \draw[setinha, cBlue] (a.north) -- (c2.south);
  }
  \only<5->{
    \node[nota=cRed, anchor=north west] (b) at ([xshift=-0.4cm,yshift=-0.3cm]mm.south) {$\overline{m} = 0$: mensagem sem nível CC};
    \draw[setinha, cRed] ([xshift=0.4cm]b.north west) -- (mm.south);
  }
}%
\onslide<6->{
\[
\boxed{\,S_i = \tikzmarknode{pc}{\frac{A^2}{2}} + \tikzmarknode{pl}{\frac{\overline{m^2}}{2}}\,}
\]}%
\par\vspace{0.6cm}%
\anotacoes{\only<7->{
  \node[nota=cOrange, anchor=north east] (c) at ([xshift=-0.3cm,yshift=-0.3cm]pc.south) {portadora: não leva informação};
  \draw[setinha, cOrange] (c.north) -- (pc.south);
  \node[nota=cGreen, anchor=north west] (d) at ([xshift=0.3cm,yshift=-0.3cm]pl.south) {bandas laterais: a mensagem};
  \draw[setinha, cGreen] (d.north) -- (pl.south);
}}%

\end{frame}

% ============================================

\begin{frame}{AM, ponto 1: ruído e SNR}

\topoAM{1}

\begin{columns}[T]
\column{0.5\textwidth}
\begin{tikzpicture}
\begin{axis}[onda, width=5.6cm, height=1.5cm, xmin=-4, xmax=4, ymin=0, ymax=1.6,
             xlabel={$f$}, ylabel={espectro no ponto 1}]
  \only<2->{
    \addplot[cRed, fill=cRed!15, draw=none, samples=2] coordinates {(-3.1,0) (-3.1,0.3) (-1.9,0.3) (-1.9,0)};
    \addplot[cRed, fill=cRed!15, draw=none, samples=2] coordinates {(1.9,0) (1.9,0.3) (3.1,0.3) (3.1,0)};
    \node[cRed, font=\tiny, anchor=south] at (axis cs:0,0.3) {$N_0/2$};
  }
  \addplot[cBlue, thick, fill=cBlue!20, samples=2] coordinates {(1.9,0) (2.5,0.8) (3.1,0)};
  \addplot[cBlue, thick, fill=cBlue!20, samples=2] coordinates {(-3.1,0) (-2.5,0.8) (-1.9,0)};
  \draw[cOrange, very thick, -{Stealth[length=2mm]}] (axis cs:2.5,0) -- (axis cs:2.5,1.45);
  \draw[cOrange, very thick, -{Stealth[length=2mm]}] (axis cs:-2.5,0) -- (axis cs:-2.5,1.45);
  \node[cOrange, font=\tiny, anchor=west] at (axis cs:2.6,1.4) {portadora};
  \node[font=\tiny, anchor=north] at (axis cs:2.5,0) {$f_c$};
  \node[font=\tiny, anchor=north] at (axis cs:-2.5,0) {$-f_c$};
\end{axis}
\end{tikzpicture}

\vspace{0.2cm}\small
\only<1>{A portadora aparece como componente de frequência em $\pm f_c$; a mensagem, nas bandas laterais.}%
\only<2->{O filtro tem a mesma banda $2W$ do DSB-SC: o ruído é o mesmo.}

\column{0.5\textwidth}
\small\compacto
\onslide<3->{
\[
\overline{n^2} = 2\times\frac{N_0}{2}\times 2W = 2N_0W
\]}%
\onslide<4->{
\[
\left(\frac{S}{N}\right)_1 = \frac{S_i}{2N_0W}
\only<5->{= \frac12\,\tikzmarknode{g1}{\frac{S_i}{N_0W}}}
\only<6->{= \frac{\gamma}{2}}
\]}%
\par\vspace{0.7cm}%
\anotacoes{\only<5->{
  \node[nota=cGreen, anchor=north] (a) at ([yshift=-0.3cm]g1.south) {$\gamma$ com a mesma $S_i$};
  \draw[setinha, cGreen] (a.north) -- (g1.south);
}}%
\onslide<7->{Igual ao DSB-SC. A diferença vai aparecer depois do detector: parte de $S_i$ é portadora.}
\end{columns}

\end{frame}

% ============================================

\begin{frame}{Envoltória com ruído}

\begin{columns}[T]
\column{0.42\textwidth}
\begin{tikzpicture}[scale=1.05, >={Stealth[length=2mm]}]
  \draw[gray!60, ->] (-0.2,0) -- (4.3,0) node[right, font=\scriptsize, gray] {fase};
  \draw[gray!60, ->] (0,-0.2) -- (0,1.5) node[above, font=\scriptsize, gray] {quadratura};
  \draw[->, cBlue, very thick] (0,0) -- (3,0) node[midway, below, font=\scriptsize] {$A + m$};
  \onslide<2->{\draw[->, cRed, very thick] (3,0) -- (3.6,0) node[below, font=\scriptsize] {$n_c$};}
  \onslide<3->{\draw[->, cOrange, very thick] (3.6,0) -- (3.6,0.55) node[right, font=\scriptsize] {$n_s$};}
  \onslide<4->{\draw[->, cGreen, very thick] (0,0) -- (3.6,0.55) node[above, font=\scriptsize, pos=0.6] {$E(t)$};}
\end{tikzpicture}

\vspace{0.2cm}\small
\only<1>{O sinal é um fasor em fase com a portadora.}%
\only<2>{$n_c$ soma na mesma direção.}%
\only<3>{$n_s$ entra a $90^\circ$.}%
\only<4->{A envoltória é o comprimento do fasor resultante.}

\column{0.58\textwidth}
\small\compacto
Entrada do detector, agrupando por portadora:
\[
r(t) = \big[A + m + n_c\big]\cos(2\pi f_c t) - n_s\sin(2\pi f_c t)
\]
\onslide<4->{
\[
E(t) = \sqrt{\big[A + m + n_c\big]^2 + \tikzmarknode{ns2}{n_s^2}}
\]}%
\par\vspace{0.7cm}%
\anotacoes{\only<5->{
  \node[nota=cOrange, anchor=north] (a) at ([xshift=-0.6cm,yshift=-0.3cm]ns2.south) {desprezível se $A + m \gg |n|$};
  \draw[setinha, cOrange] (a.north) -- (ns2.south);
}}%
\onslide<6->{Ruído fraco frente à portadora:
\[
\boxed{\,E(t) \approx A + m(t) + n_c(t)\,}
\]}%
\onslide<7->{O detector de envoltória entrega a componente em fase, como o detector coerente.}
\end{columns}

\end{frame}

% ============================================

\begin{frame}{AM, ponto 2: sinal e ruído}

\small\compacto
\[
E(t) \approx \tikzmarknode{eA}{A} + m(t) + n_c(t)
\]
\anotacoes{
  \node[nota=cOrange, anchor=east] (a) at ([xshift=-1.2cm]eA.west) {CC: o bloqueio retira};
  \draw[setinha, cOrange] (a.east) -- (eA.west);
}%
\onslide<2->{
\[
y(t) = \tikzmarknode{ym}{m(t)} + \tikzmarknode{yn}{n_c(t)}
\]}%
\anotacoes{\only<2->{
  \node[nota=cBlue, anchor=east] (b) at ([xshift=-0.5cm]ym.west) {sinal};
  \draw[setinha, cBlue] (b.east) -- (ym.west);
  \node[nota=cRed, anchor=west] (c) at ([xshift=0.5cm]yn.east) {ruído};
  \draw[setinha, cRed] (c.west) -- (yn.east);
}}%
\onslide<3->{Sem o fator $\tfrac12$ do misturador:}
\onslide<3->{
\[
S_o = \overline{m^2}, \qquad N_o = \tikzmarknode{nc2}{\overline{n_c^2}}
\only<4->{= 2N_0W}
\]}%
\par\vspace{0.5cm}%
\anotacoes{\only<4->{
  \node[nota=cGreen, anchor=north west] (d) at ([xshift=0.2cm,yshift=-0.25cm]nc2.south) {slide de fase e quadratura};
  \draw[setinha, cGreen] ([xshift=0.3cm]d.north west) -- (nc2.south);
}}%
\onslide<5->{
\[
\left(\frac{S}{N}\right)_2 = \frac{\overline{m^2}}{2N_0W}
\]}

\end{frame}

% ============================================

\begin{frame}{AM, ponto 2: SNR em função de $\gamma$}

\small\compacto
Para comparar com a banda base, a SNR precisa aparecer com $\gamma = S_i/(N_0W)$.
\onslide<2->{Multiplicando e dividindo por $A^2 + \overline{m^2}$:}
\onslide<2->{
\[
\left(\frac{S}{N}\right)_2 = \frac{\overline{m^2}}{2N_0W}
= \tikzmarknode{eta}{\frac{\overline{m^2}}{A^2 + \overline{m^2}}}\cdot
  \tikzmarknode{gam}{\frac{A^2 + \overline{m^2}}{2N_0W}}
\]}%
\par\vspace{1.6cm}%
\anotacoes{
  \only<3->{
    \node[nota=cGreen, anchor=north east] (a) at ([xshift=1.2cm,yshift=-0.3cm]gam.south) {$\dfrac{A^2+\overline{m^2}}{2} = S_i$: vale $\dfrac{S_i}{N_0W} = \gamma$};
    \draw[setinha, cGreen] ([xshift=-1.2cm]a.north east) -- (gam.south);
  }
  \only<4->{
    \node[nota=cOrange, anchor=north east] (b) at ([xshift=-0.3cm,yshift=-1.3cm]eta.south) {$\eta$: fração de $S_i$ nas bandas laterais};
    \draw[setinha, cOrange] (b.north) -- (eta.south);
  }
}%
\onslide<5->{
\[
\boxed{\,\left(\frac{S}{N}\right)_2 = \eta\,\gamma, \qquad \eta = \frac{\overline{m^2}}{A^2 + \overline{m^2}} < 1\,}
\]}%
\onslide<6->{Sempre pior que o DSB-SC: a portadora conta em $S_i$ e não vira sinal.}
\onslide<7->{Ganho de detecção: $\eta\gamma/(\gamma/2) = 2\eta$.}

\end{frame}

% ============================================

\begin{frame}{AM com modulação por tom}

\begin{columns}[T]
\column{0.55\textwidth}
\small\compacto
Tom com índice $\mu$:
\[
m(t) = \mu A\cos(2\pi f_m t)
\]
\onslide<2->{
\[
\overline{m^2} = \mu^2A^2\,\tikzmarknode{cm}{\overline{\cos^2(2\pi f_m t)}}
\only<3->{= \frac{\mu^2A^2}{2}}
\]}%
\par\vspace{0.4cm}%
\anotacoes{\only<3->{
  \node[nota=cBlue, anchor=north] (a) at ([yshift=-0.2cm]cm.south) {média $\tfrac12$};
  \draw[setinha, cBlue] (a.north) -- (cm.south);
}}%
\onslide<4->{
\[
\eta = \frac{\mu^2A^2/2}{A^2 + \mu^2A^2/2}
\only<5->{= \tikzmarknode{et}{\frac{\mu^2}{2 + \mu^2}}}
\]}%
\anotacoes{\only<5->{
  \node[nota=cGreen, anchor=west] (b) at ([xshift=1.0cm]et.east) {$A^2$ cancela};
  \draw[setinha, cGreen] (b.west) -- (et.east);
}}%
\onslide<7->{Com $\mu = 1$, o melhor caso sem sobremodulação:
\[
\eta = \frac13 \;\Rightarrow\; \left(\frac{S}{N}\right)_2 = \frac{\gamma}{3}, \quad 4{,}77\ \text{dB abaixo}
\]}

\column{0.45\textwidth}
\onslide<6->{
\begin{tikzpicture}
\begin{axis}[espec, width=4.6cm, height=3.6cm, xmin=0, xmax=1.05, ymax=0.4,
             xlabel={$\mu$}, ylabel={$\eta$}, xtick={0,0.5,1}, ytick={0,0.1,0.2,1/3},
             yticklabels={$0$,$0{,}1$,$0{,}2$,$\tfrac13$}, grid=major]
  \addplot[cBlue, thick, domain=0:1, samples=60] {x^2/(2+x^2)};
  \only<7->{\addplot[cRed, only marks, mark=*, mark size=2pt] coordinates {(1,1/3)};}
\end{axis}
\end{tikzpicture}}
\end{columns}

\end{frame}

% ============================================

\begin{frame}{AM: limiar com ruído forte}

\begin{columns}[T]
\column{0.4\textwidth}
\begin{tikzpicture}[scale=0.95, >={Stealth[length=2mm]}]
  \draw[gray!60, ->] (-0.2,0) -- (3.2,0);
  \draw[gray!60, ->] (0,-0.2) -- (0,2.9);
  \draw[->, cBlue, very thick] (0,0) -- (0.7,0) node[below, font=\scriptsize] {$A + m$};
  \onslide<2->{
    \draw[->, cRed, very thick] (0.7,0) -- ++(65:2.8) node[right, font=\scriptsize] {$R_n$};
    \draw[cRed] (1.2,0) arc (0:65:0.5);
    \node[cRed, font=\scriptsize] at (1.4,0.35) {$\theta_n$};
  }
  \onslide<3->{\draw[->, cGreen, very thick] (0,0) -- ($(0.7,0)+(65:2.8)$) node[midway, left, font=\scriptsize] {$E$};}
\end{tikzpicture}

\vspace{0.1cm}\small
\only<1>{Agora o ruído é maior que a portadora.}%
\only<2>{Ruído em forma polar: $n_c = R_n\cos\theta_n$, $n_s = R_n\sin\theta_n$.}%
\only<3->{A envoltória segue o ruído, que gira com fase aleatória.}

\onslide<8->{\textcolor{cRed}{Limiar do detector de envoltória.} O detector coerente não tem limiar e dá $\eta\gamma$ com qualquer SNR.}

\column{0.6\textwidth}
\small\compacto
\onslide<2->{
\[
E^2 = \big[A + m + R_n\cos\theta_n\big]^2 + R_n^2\sin^2\theta_n
\]}%
\onslide<3->{
\[
E^2 = \tikzmarknode{x2}{(A+m)^2} + 2(A+m)R_n\cos\theta_n + R_n^2
\]}%
\par\vspace{0.6cm}%
\anotacoes{\only<4->{
  \node[nota=cOrange, anchor=north] (a) at ([yshift=-0.25cm]x2.south) {desprezível se $R_n \gg A + m$};
  \draw[setinha, cOrange] (a.north) -- (x2.south);
}}%
\onslide<5->{
\[
E \approx R_n\,\tikzmarknode{sq}{\sqrt{1 + \frac{2(A+m)\cos\theta_n}{R_n}}}
\only<6->{\approx R_n + \big(A + m\big)\cos\theta_n}
\]}%
\par\vspace{0.6cm}%
\anotacoes{\only<6->{
  \node[nota=cBlue, anchor=north] (b) at ([yshift=-0.25cm]sq.south) {$\sqrt{1+x} \approx 1 + x/2$};
  \draw[setinha, cBlue] (b.north) -- (sq.south);
}}%
\onslide<7->{A mensagem só aparece como $m\cos\theta_n$, multiplicada por um ruído: não há parcela proporcional a $m(t)$.}
\end{columns}

\end{frame}

% ============================================
